% Options for packages loaded elsewhere
\PassOptionsToPackage{unicode}{hyperref}
\PassOptionsToPackage{hyphens}{url}
%
\documentclass[
  11pt,
]{article}
\usepackage{amsmath,amssymb}
\usepackage{iftex}
\ifPDFTeX
  \usepackage[T1]{fontenc}
  \usepackage[utf8]{inputenc}
  \usepackage{textcomp} % provide euro and other symbols
\else % if luatex or xetex
  \usepackage{unicode-math} % this also loads fontspec
  \defaultfontfeatures{Scale=MatchLowercase}
  \defaultfontfeatures[\rmfamily]{Ligatures=TeX,Scale=1}
\fi
\usepackage{lmodern}
\ifPDFTeX\else
  % xetex/luatex font selection
\fi
% Use upquote if available, for straight quotes in verbatim environments
\IfFileExists{upquote.sty}{\usepackage{upquote}}{}
\IfFileExists{microtype.sty}{% use microtype if available
  \usepackage[]{microtype}
  \UseMicrotypeSet[protrusion]{basicmath} % disable protrusion for tt fonts
}{}
\makeatletter
\@ifundefined{KOMAClassName}{% if non-KOMA class
  \IfFileExists{parskip.sty}{%
    \usepackage{parskip}
  }{% else
    \setlength{\parindent}{0pt}
    \setlength{\parskip}{6pt plus 2pt minus 1pt}}
}{% if KOMA class
  \KOMAoptions{parskip=half}}
\makeatother
\usepackage{xcolor}
\usepackage[margin=2.2cm]{geometry}
\usepackage{longtable,booktabs,array}
\usepackage{calc} % for calculating minipage widths
% Correct order of tables after \paragraph or \subparagraph
\usepackage{etoolbox}
\makeatletter
\patchcmd\longtable{\par}{\if@noskipsec\mbox{}\fi\par}{}{}
\makeatother
% Allow footnotes in longtable head/foot
\IfFileExists{footnotehyper.sty}{\usepackage{footnotehyper}}{\usepackage{footnote}}
\makesavenoteenv{longtable}
\setlength{\emergencystretch}{3em} % prevent overfull lines
\providecommand{\tightlist}{%
  \setlength{\itemsep}{0pt}\setlength{\parskip}{0pt}}
\setcounter{secnumdepth}{-\maxdimen} % remove section numbering
\ifLuaTeX
  \usepackage{selnolig}  % disable illegal ligatures
\fi
\IfFileExists{bookmark.sty}{\usepackage{bookmark}}{\usepackage{hyperref}}
\IfFileExists{xurl.sty}{\usepackage{xurl}}{} % add URL line breaks if available
\urlstyle{same}
\hypersetup{
  pdftitle={Worksheet 6A: Statistics and Probability},
  hidelinks,
  pdfcreator={LaTeX via pandoc}}

\title{Worksheet 6A: Statistics and Probability}
\author{}
\date{}

\begin{document}
\maketitle

\begin{quote}
\textbf{What This Worksheet Is About}

Statistics and probability are the mathematical foundation of machine
learning. Every time an AI makes a prediction, classifies an image, or
recommends a video, it's using these concepts.

This worksheet connects the basics you've learned to real applications
--- from training neural networks to analyzing scientific data.
\end{quote}

\begin{center}\rule{0.5\linewidth}{0.5pt}\end{center}

\hypertarget{part-a-measures-of-central-tendency-mean-median-mode-18-problems}{%
\subsection{Part A: Measures of Central Tendency --- Mean, Median, Mode
(18
problems)}\label{part-a-measures-of-central-tendency-mean-median-mode-18-problems}}

\hypertarget{the-formulas}{%
\subsubsection{The Formulas}\label{the-formulas}}

\textbf{Mean (Average):}
\(\bar{x} = \frac{\sum x_i}{n} = \frac{x_1 + x_2 + \cdots + x_n}{n}\)

\textbf{Median:} The middle value when data is sorted (or average of two
middle values if \(n\) is even)

\textbf{Mode:} The most frequent value

\begin{center}\rule{0.5\linewidth}{0.5pt}\end{center}

\textbf{Calculate the mean, median, and mode:}

\textbf{1.} Test scores: 78, 82, 85, 85, 88, 90, 92

\textbf{2.} Response times (ms): 120, 135, 140, 145, 145, 150, 180

\textbf{3.} Daily website visitors: 1200, 1350, 1400, 1450, 1500, 1550,
8500

\begin{itemize}
\item
  \begin{enumerate}
  \def\labelenumi{(\alph{enumi})}
  \tightlist
  \item
    Calculate all three measures.
  \end{enumerate}
\item
  \begin{enumerate}
  \def\labelenumi{(\alph{enumi})}
  \setcounter{enumi}{1}
  \tightlist
  \item
    Which measure is most affected by the outlier (8500)?
  \end{enumerate}
\item
  \begin{enumerate}
  \def\labelenumi{(\alph{enumi})}
  \setcounter{enumi}{2}
  \tightlist
  \item
    Which measure best represents a ``typical'' day?
  \end{enumerate}
\end{itemize}

\begin{center}\rule{0.5\linewidth}{0.5pt}\end{center}

** AI Connection: Training Data Quality**

\textbf{4.} A dataset of house prices (in €1000s) contains: 250, 275,
290, 310, 325, 340, 2500

The last value is a mansion that got mixed in with regular houses.

\begin{itemize}
\item
  \begin{enumerate}
  \def\labelenumi{(\alph{enumi})}
  \tightlist
  \item
    Calculate the mean with and without the outlier.
  \end{enumerate}
\item
  \begin{enumerate}
  \def\labelenumi{(\alph{enumi})}
  \setcounter{enumi}{1}
  \tightlist
  \item
    Calculate the median with and without the outlier.
  \end{enumerate}
\item
  \begin{enumerate}
  \def\labelenumi{(\alph{enumi})}
  \setcounter{enumi}{2}
  \tightlist
  \item
    Why might a data scientist remove outliers before training a model?
  \end{enumerate}
\end{itemize}

\textbf{5.} An image classification dataset has these label counts:

\begin{longtable}[]{@{}ll@{}}
\toprule\noalign{}
Category & Count \\
\midrule\noalign{}
\endhead
\bottomrule\noalign{}
\endlastfoot
Cat & 5000 \\
Dog & 4800 \\
Bird & 1200 \\
Fish & 500 \\
\end{longtable}

\begin{itemize}
\item
  \begin{enumerate}
  \def\labelenumi{(\alph{enumi})}
  \tightlist
  \item
    What is the mean number of images per category?
  \end{enumerate}
\item
  \begin{enumerate}
  \def\labelenumi{(\alph{enumi})}
  \setcounter{enumi}{1}
  \tightlist
  \item
    What is the median?
  \end{enumerate}
\item
  \begin{enumerate}
  \def\labelenumi{(\alph{enumi})}
  \setcounter{enumi}{2}
  \tightlist
  \item
    What is the mode?
  \end{enumerate}
\item
  \begin{enumerate}
  \def\labelenumi{(\alph{enumi})}
  \setcounter{enumi}{3}
  \tightlist
  \item
    This is called an ``imbalanced dataset.'' Why might this be a
    problem for training?
  \end{enumerate}
\end{itemize}

\begin{center}\rule{0.5\linewidth}{0.5pt}\end{center}

** Science Connection: Experimental Measurements**

\textbf{6.} A physicist measures the speed of light in 5 trials (in
units of \(10^8\) m/s):

2.998, 2.995, 3.001, 2.997, 2.999

\begin{itemize}
\item
  \begin{enumerate}
  \def\labelenumi{(\alph{enumi})}
  \tightlist
  \item
    Calculate the mean.
  \end{enumerate}
\item
  \begin{enumerate}
  \def\labelenumi{(\alph{enumi})}
  \setcounter{enumi}{1}
  \tightlist
  \item
    The accepted value is 2.998. How close is the experimental mean?
  \end{enumerate}
\end{itemize}

\textbf{7.} A chemist measures the pH of a solution 7 times:

7.02, 6.98, 7.01, 6.99, 7.00, 7.01, 7.00

\begin{itemize}
\item
  \begin{enumerate}
  \def\labelenumi{(\alph{enumi})}
  \tightlist
  \item
    Find mean, median, and mode.
  \end{enumerate}
\item
  \begin{enumerate}
  \def\labelenumi{(\alph{enumi})}
  \setcounter{enumi}{1}
  \tightlist
  \item
    Are they all similar? What does this suggest about the measurement
    consistency?
  \end{enumerate}
\end{itemize}

\begin{center}\rule{0.5\linewidth}{0.5pt}\end{center}

** Game Analytics**

\textbf{8.} Player scores in a mobile game level: 1250, 1480, 1520,
1550, 1600, 1650, 1700, 1800, 1850, 2400

\begin{itemize}
\item
  \begin{enumerate}
  \def\labelenumi{(\alph{enumi})}
  \tightlist
  \item
    Find the mean and median.
  \end{enumerate}
\item
  \begin{enumerate}
  \def\labelenumi{(\alph{enumi})}
  \setcounter{enumi}{1}
  \tightlist
  \item
    If the game designer wants to set the ``3-star threshold'' at a
    score that about half of players achieve, should they use the mean
    or median?
  \end{enumerate}
\end{itemize}

\textbf{9.} Daily active users for an app over two weeks:

Week 1: 10000, 12000, 15000, 14000, 13000, 8000, 7000 Week 2: 11000,
13000, 16000, 15000, 14000, 9000, 8000

\begin{itemize}
\item
  \begin{enumerate}
  \def\labelenumi{(\alph{enumi})}
  \tightlist
  \item
    Calculate the mean for each week.
  \end{enumerate}
\item
  \begin{enumerate}
  \def\labelenumi{(\alph{enumi})}
  \setcounter{enumi}{1}
  \tightlist
  \item
    Is the app growing? By what percentage?
  \end{enumerate}
\end{itemize}

\begin{center}\rule{0.5\linewidth}{0.5pt}\end{center}

\textbf{Weighted Mean:}

\textbf{10.} A student's grade is calculated as: - Homework: 25\%
weight, score 85 - Midterm: 30\% weight, score 78\\
- Final: 45\% weight, score 88

Calculate the weighted mean (final grade):
\[\bar{x}_w = \frac{\sum w_i x_i}{\sum w_i}\]

\textbf{11.} In a neural network, outputs from three layers are combined
with weights: - Layer 1: output 0.7, weight 0.5 - Layer 2: output 0.3,
weight 0.3 - Layer 3: output 0.8, weight 0.2

Calculate the weighted average output.

\begin{center}\rule{0.5\linewidth}{0.5pt}\end{center}

\hypertarget{part-b-measures-of-spread-range-variance-standard-deviation-16-problems}{%
\subsection{Part B: Measures of Spread --- Range, Variance, Standard
Deviation (16
problems)}\label{part-b-measures-of-spread-range-variance-standard-deviation-16-problems}}

\hypertarget{the-formulas-1}{%
\subsubsection{The Formulas}\label{the-formulas-1}}

\textbf{Range:} \(R = x_{max} - x_{min}\)

\textbf{Variance:} \(\sigma^2 = \frac{\sum (x_i - \bar{x})^2}{n}\)

\textbf{Standard Deviation:} \(\sigma = \sqrt{\sigma^2}\)

\emph{Note: For samples, divide by \((n-1)\) instead of \(n\) --- this
is called Bessel's correction.}

\begin{center}\rule{0.5\linewidth}{0.5pt}\end{center}

\textbf{Calculate the range, variance, and standard deviation:}

\textbf{12.} Data: 4, 6, 8, 10, 12

\begin{itemize}
\item
  \begin{enumerate}
  \def\labelenumi{(\alph{enumi})}
  \tightlist
  \item
    Find the mean.
  \end{enumerate}
\item
  \begin{enumerate}
  \def\labelenumi{(\alph{enumi})}
  \setcounter{enumi}{1}
  \tightlist
  \item
    Find \((x_i - \bar{x})\) for each value.
  \end{enumerate}
\item
  \begin{enumerate}
  \def\labelenumi{(\alph{enumi})}
  \setcounter{enumi}{2}
  \tightlist
  \item
    Find \((x_i - \bar{x})^2\) for each value.
  \end{enumerate}
\item
  \begin{enumerate}
  \def\labelenumi{(\alph{enumi})}
  \setcounter{enumi}{3}
  \tightlist
  \item
    Find the variance (divide by \(n = 5\)).
  \end{enumerate}
\item
  \begin{enumerate}
  \def\labelenumi{(\alph{enumi})}
  \setcounter{enumi}{4}
  \tightlist
  \item
    Find the standard deviation.
  \end{enumerate}
\end{itemize}

\textbf{13.} Data: 7, 7, 8, 8, 8, 9, 9

\textbf{14.} Data: 2, 5, 8, 11, 14

\begin{center}\rule{0.5\linewidth}{0.5pt}\end{center}

** AI Connection: Model Consistency**

\textbf{15.} Two machine learning models are tested 5 times each on
accuracy (\%):

Model A: 85, 87, 84, 86, 88 Model B: 70, 95, 75, 90, 80

\begin{itemize}
\item
  \begin{enumerate}
  \def\labelenumi{(\alph{enumi})}
  \tightlist
  \item
    Calculate the mean accuracy for each.
  \end{enumerate}
\item
  \begin{enumerate}
  \def\labelenumi{(\alph{enumi})}
  \setcounter{enumi}{1}
  \tightlist
  \item
    Calculate the standard deviation for each.
  \end{enumerate}
\item
  \begin{enumerate}
  \def\labelenumi{(\alph{enumi})}
  \setcounter{enumi}{2}
  \tightlist
  \item
    Which model is more consistent?
  \end{enumerate}
\item
  \begin{enumerate}
  \def\labelenumi{(\alph{enumi})}
  \setcounter{enumi}{3}
  \tightlist
  \item
    In production, would you prefer a consistent model or one with
    higher peaks?
  \end{enumerate}
\end{itemize}

\textbf{16.} A chatbot's response times (seconds) over 6 queries:

Version 1.0: 1.2, 1.3, 1.1, 1.4, 1.2, 1.3 Version 2.0: 0.8, 2.1, 0.7,
1.8, 0.9, 1.7

\begin{itemize}
\item
  \begin{enumerate}
  \def\labelenumi{(\alph{enumi})}
  \tightlist
  \item
    Which version has lower mean response time?
  \end{enumerate}
\item
  \begin{enumerate}
  \def\labelenumi{(\alph{enumi})}
  \setcounter{enumi}{1}
  \tightlist
  \item
    Which version is more consistent (lower standard deviation)?
  \end{enumerate}
\item
  \begin{enumerate}
  \def\labelenumi{(\alph{enumi})}
  \setcounter{enumi}{2}
  \tightlist
  \item
    Users hate unpredictable wait times. Which version would you deploy?
  \end{enumerate}
\end{itemize}

\begin{center}\rule{0.5\linewidth}{0.5pt}\end{center}

** Science Connection: Measurement Precision**

\textbf{17.} Two instruments measure the same voltage (in volts):

Instrument A: 5.02, 4.98, 5.01, 4.99, 5.00 Instrument B: 5.10, 4.90,
5.15, 4.85, 5.00

\begin{itemize}
\item
  \begin{enumerate}
  \def\labelenumi{(\alph{enumi})}
  \tightlist
  \item
    Both have the same mean. Verify this.
  \end{enumerate}
\item
  \begin{enumerate}
  \def\labelenumi{(\alph{enumi})}
  \setcounter{enumi}{1}
  \tightlist
  \item
    Calculate the standard deviation for each.
  \end{enumerate}
\item
  \begin{enumerate}
  \def\labelenumi{(\alph{enumi})}
  \setcounter{enumi}{2}
  \tightlist
  \item
    Which instrument is more precise?
  \end{enumerate}
\end{itemize}

\textbf{18.} In manufacturing, parts must be within tolerance. A machine
produces bolts with target diameter 10.00 mm. A sample of 5 bolts
measures:

10.02, 9.98, 10.01, 9.99, 10.00

\begin{itemize}
\item
  \begin{enumerate}
  \def\labelenumi{(\alph{enumi})}
  \tightlist
  \item
    Find the mean and standard deviation.
  \end{enumerate}
\item
  \begin{enumerate}
  \def\labelenumi{(\alph{enumi})}
  \setcounter{enumi}{1}
  \tightlist
  \item
    If tolerance is ±0.05 mm (so 9.95 to 10.05 is acceptable), are all
    bolts within tolerance?
  \end{enumerate}
\item
  \begin{enumerate}
  \def\labelenumi{(\alph{enumi})}
  \setcounter{enumi}{2}
  \tightlist
  \item
    If the standard deviation were 0.03 mm instead, what percentage of
    bolts would you expect to be outside tolerance? (Assume normal
    distribution where \textasciitilde68\% are within 1 SD of mean)
  \end{enumerate}
\end{itemize}

\begin{center}\rule{0.5\linewidth}{0.5pt}\end{center}

** Understanding Standard Deviation:**

\textbf{19.} For a normal distribution: - \textasciitilde68\% of data
falls within 1 standard deviation of the mean - \textasciitilde95\%
falls within 2 standard deviations - \textasciitilde99.7\% falls within
3 standard deviations

IQ scores have mean 100 and standard deviation 15.

\begin{itemize}
\item
  \begin{enumerate}
  \def\labelenumi{(\alph{enumi})}
  \tightlist
  \item
    What range contains about 68\% of people?
  \end{enumerate}
\item
  \begin{enumerate}
  \def\labelenumi{(\alph{enumi})}
  \setcounter{enumi}{1}
  \tightlist
  \item
    What range contains about 95\% of people?
  \end{enumerate}
\item
  \begin{enumerate}
  \def\labelenumi{(\alph{enumi})}
  \setcounter{enumi}{2}
  \tightlist
  \item
    What percentage of people have IQ above 130?
  \end{enumerate}
\end{itemize}

\textbf{20.} Human height (adult males) has mean ≈ 175 cm and SD ≈ 7 cm.

\begin{itemize}
\item
  \begin{enumerate}
  \def\labelenumi{(\alph{enumi})}
  \tightlist
  \item
    What range contains the middle 68\%?
  \end{enumerate}
\item
  \begin{enumerate}
  \def\labelenumi{(\alph{enumi})}
  \setcounter{enumi}{1}
  \tightlist
  \item
    How tall would someone need to be to be in the top 2.5\%?
  \end{enumerate}
\item
  \begin{enumerate}
  \def\labelenumi{(\alph{enumi})}
  \setcounter{enumi}{2}
  \tightlist
  \item
    A basketball team requires players at least 190 cm tall.
    Approximately what percentage of men qualify?
  \end{enumerate}
\end{itemize}

\begin{center}\rule{0.5\linewidth}{0.5pt}\end{center}

\hypertarget{part-c-basic-probability-20-problems}{%
\subsection{Part C: Basic Probability (20
problems)}\label{part-c-basic-probability-20-problems}}

\hypertarget{the-fundamentals}{%
\subsubsection{The Fundamentals}\label{the-fundamentals}}

\[P(\text{event}) = \frac{\text{Number of favorable outcomes}}{\text{Total possible outcomes}}\]

\textbf{Key Rules:} - \(0 \leq P(A) \leq 1\) -
\(P(\text{certain event}) = 1\) - \(P(\text{impossible event}) = 0\) -
\(P(\text{not } A) = 1 - P(A)\)

\begin{center}\rule{0.5\linewidth}{0.5pt}\end{center}

\textbf{21.} A standard deck has 52 cards (4 suits × 13 values). Find: -
(a) \(P(\text{drawing a heart})\) - (b) \(P(\text{drawing an ace})\) -
(c) \(P(\text{drawing the ace of hearts})\) - (d)
\(P(\text{not drawing a face card})\) (face cards = J, Q, K)

\textbf{22.} A bag contains 5 red, 3 blue, and 2 green marbles. Find: -
(a) \(P(\text{red})\) - (b) \(P(\text{not red})\) - (c)
\(P(\text{red or blue})\)

\begin{center}\rule{0.5\linewidth}{0.5pt}\end{center}

** AI Connection: Classification Probabilities**

\textbf{23.} A spam filter classifies 1000 emails: - 800 are correctly
identified (700 real emails marked ``not spam'', 100 spam marked
``spam'') - 200 are incorrectly identified (150 real emails marked
``spam'', 50 spam marked ``not spam'')

\begin{itemize}
\item
  \begin{enumerate}
  \def\labelenumi{(\alph{enumi})}
  \tightlist
  \item
    What is \(P(\text{correct classification})\)?
  \end{enumerate}
\item
  \begin{enumerate}
  \def\labelenumi{(\alph{enumi})}
  \setcounter{enumi}{1}
  \tightlist
  \item
    What is \(P(\text{false positive})\) --- real email marked as spam?
  \end{enumerate}
\item
  \begin{enumerate}
  \def\labelenumi{(\alph{enumi})}
  \setcounter{enumi}{2}
  \tightlist
  \item
    What is \(P(\text{false negative})\) --- spam marked as real?
  \end{enumerate}
\item
  \begin{enumerate}
  \def\labelenumi{(\alph{enumi})}
  \setcounter{enumi}{3}
  \tightlist
  \item
    Which error is worse for the user?
  \end{enumerate}
\end{itemize}

\textbf{24.} An image classifier sees 500 images: - 200 cats (correctly
identifies 180) - 200 dogs (correctly identifies 170) - 100 birds
(correctly identifies 80)

\begin{itemize}
\item
  \begin{enumerate}
  \def\labelenumi{(\alph{enumi})}
  \tightlist
  \item
    What is \(P(\text{correct | cat})\) --- probability of correct given
    it's a cat?
  \end{enumerate}
\item
  \begin{enumerate}
  \def\labelenumi{(\alph{enumi})}
  \setcounter{enumi}{1}
  \tightlist
  \item
    What is \(P(\text{correct | dog})\)?
  \end{enumerate}
\item
  \begin{enumerate}
  \def\labelenumi{(\alph{enumi})}
  \setcounter{enumi}{2}
  \tightlist
  \item
    What is \(P(\text{correct | bird})\)?
  \end{enumerate}
\item
  \begin{enumerate}
  \def\labelenumi{(\alph{enumi})}
  \setcounter{enumi}{3}
  \tightlist
  \item
    What is the overall accuracy?
  \end{enumerate}
\end{itemize}

\begin{center}\rule{0.5\linewidth}{0.5pt}\end{center}

** Science Connection: Medical Testing**

\textbf{25.} A disease affects 1\% of the population. A test for the
disease: - Correctly identifies 95\% of sick people (sensitivity) -
Correctly identifies 90\% of healthy people (specificity)

In a population of 10,000 people: - (a) How many are actually sick?
Healthy? - (b) Of the sick, how many test positive? - (c) Of the
healthy, how many test positive (false positives)? - (d) If someone
tests positive, what's the probability they're actually sick?

\emph{This is the famous ``base rate fallacy'' --- the answer may
surprise you!}

\begin{center}\rule{0.5\linewidth}{0.5pt}\end{center}

** Game Probability**

\textbf{26.} A loot box in a game contains: - Common item: 70\% chance -
Rare item: 20\% chance - Epic item: 8\% chance - Legendary item: 2\%
chance

\begin{itemize}
\item
  \begin{enumerate}
  \def\labelenumi{(\alph{enumi})}
  \tightlist
  \item
    What is \(P(\text{rare or better})\)?
  \end{enumerate}
\item
  \begin{enumerate}
  \def\labelenumi{(\alph{enumi})}
  \setcounter{enumi}{1}
  \tightlist
  \item
    If you open 3 boxes, what's the probability of getting no legendary
    items?
  \end{enumerate}
\item
  \begin{enumerate}
  \def\labelenumi{(\alph{enumi})}
  \setcounter{enumi}{2}
  \tightlist
  \item
    A player opens 50 boxes and gets no legendary items. How ``unlucky''
    is this? Calculate
    \(P(\text{no legendary in 50 boxes}) = 0.98^{50}\)
  \end{enumerate}
\end{itemize}

\textbf{27.} In a card game, you need to draw 2 specific cards from a
20-card deck: - (a) What is \(P(\text{first card is correct})\)? - (b)
Given the first card was correct, what is
\(P(\text{second card is correct})\)? - (c) What is
\(P(\text{both cards correct})\)?

\begin{center}\rule{0.5\linewidth}{0.5pt}\end{center}

\hypertarget{addition-rule-or}{%
\subsubsection{Addition Rule (OR)}\label{addition-rule-or}}

\[P(A \text{ or } B) = P(A) + P(B) - P(A \text{ and } B)\]

For mutually exclusive events: \(P(A \text{ or } B) = P(A) + P(B)\)

\textbf{28.} Rolling a die: - (a) \(P(\text{even or greater than 4})\) -
(b) Are ``even'' and ``greater than 4'' mutually exclusive?

\textbf{29.} Drawing from a deck: - (a) \(P(\text{heart or face card})\)
- (b) Why do we subtract \(P(\text{heart and face card})\)?

\begin{center}\rule{0.5\linewidth}{0.5pt}\end{center}

\hypertarget{multiplication-rule-and}{%
\subsubsection{Multiplication Rule
(AND)}\label{multiplication-rule-and}}

\textbf{Independent events:} \(P(A \text{ and } B) = P(A) \times P(B)\)

\textbf{Dependent events:} \(P(A \text{ and } B) = P(A) \times P(B|A)\)

\textbf{30.} Flipping a coin and rolling a die: - (a)
\(P(\text{heads and 6})\) - (b) \(P(\text{heads and even})\)

\textbf{31.} Drawing two cards without replacement: - (a)
\(P(\text{both aces})\) - (b) \(P(\text{both hearts})\)

\begin{center}\rule{0.5\linewidth}{0.5pt}\end{center}

** AI Connection: Independent Events in ML**

\textbf{32.} A neural network has 3 layers, each with 95\% chance of
working correctly (independently).

\begin{itemize}
\item
  \begin{enumerate}
  \def\labelenumi{(\alph{enumi})}
  \tightlist
  \item
    What is \(P(\text{all layers work})\)?
  \end{enumerate}
\item
  \begin{enumerate}
  \def\labelenumi{(\alph{enumi})}
  \setcounter{enumi}{1}
  \tightlist
  \item
    What is \(P(\text{at least one layer fails})\)?
  \end{enumerate}
\item
  \begin{enumerate}
  \def\labelenumi{(\alph{enumi})}
  \setcounter{enumi}{2}
  \tightlist
  \item
    If the network has 10 layers, each 99\% reliable, what is
    \(P(\text{all work})\)?
  \end{enumerate}
\end{itemize}

\textbf{33.} A self-driving car uses 4 independent sensors: - Camera:
99\% reliable - LIDAR: 98\% reliable - Radar: 97\% reliable -
Ultrasonic: 95\% reliable

\begin{itemize}
\item
  \begin{enumerate}
  \def\labelenumi{(\alph{enumi})}
  \tightlist
  \item
    What is \(P(\text{all sensors work})\)?
  \end{enumerate}
\item
  \begin{enumerate}
  \def\labelenumi{(\alph{enumi})}
  \setcounter{enumi}{1}
  \tightlist
  \item
    The car needs at least one sensor working. What is
    \(P(\text{at least one works})\)?
  \end{enumerate}
\end{itemize}

\begin{center}\rule{0.5\linewidth}{0.5pt}\end{center}

\hypertarget{part-d-counting-principles-16-problems}{%
\subsection{Part D: Counting Principles (16
problems)}\label{part-d-counting-principles-16-problems}}

\hypertarget{fundamental-counting-principle}{%
\subsubsection{Fundamental Counting
Principle}\label{fundamental-counting-principle}}

If task 1 can be done in \(m\) ways and task 2 can be done in \(n\)
ways, then both tasks can be done in \(m \times n\) ways.

\begin{center}\rule{0.5\linewidth}{0.5pt}\end{center}

\textbf{34.} A password has 4 characters: - First character: letter (26
options) - Second: digit (10 options) - Third: letter (26 options) -
Fourth: digit (10 options)

How many possible passwords are there?

\textbf{35.} A restaurant offers: - 4 appetizers - 6 main courses - 3
desserts

How many different 3-course meals are possible?

\begin{center}\rule{0.5\linewidth}{0.5pt}\end{center}

\hypertarget{permutations-order-matters}{%
\subsubsection{Permutations (Order
Matters)}\label{permutations-order-matters}}

\[P(n, r) = \frac{n!}{(n-r)!}\]

Number of ways to arrange \(r\) items from \(n\) items.

\textbf{36.} How many ways can 5 people line up for a photo?

\textbf{37.} How many 3-letter ``words'' (any letters, no repeats) can
be made from A, B, C, D, E?

\textbf{38.} In a race with 8 runners, how many ways can gold, silver,
and bronze be awarded?

\begin{center}\rule{0.5\linewidth}{0.5pt}\end{center}

\hypertarget{combinations-order-doesnt-matter}{%
\subsubsection{Combinations (Order Doesn't
Matter)}\label{combinations-order-doesnt-matter}}

\[C(n, r) = \binom{n}{r} = \frac{n!}{r!(n-r)!}\]

Number of ways to choose \(r\) items from \(n\) items.

\textbf{39.} How many ways can you choose 3 students from a class of 10
to form a committee?

\textbf{40.} How many 5-card poker hands are possible from a 52-card
deck?

\textbf{41.} A lottery requires choosing 6 numbers from 1-49. How many
possible tickets?

\begin{center}\rule{0.5\linewidth}{0.5pt}\end{center}

** AI Connection: Feature Selection**

\textbf{42.} A data scientist has 10 features (variables) and wants to
train models using different subsets:

\begin{itemize}
\item
  \begin{enumerate}
  \def\labelenumi{(\alph{enumi})}
  \tightlist
  \item
    How many ways to choose exactly 3 features?
  \end{enumerate}
\item
  \begin{enumerate}
  \def\labelenumi{(\alph{enumi})}
  \setcounter{enumi}{1}
  \tightlist
  \item
    How many ways to choose exactly 5 features?
  \end{enumerate}
\item
  \begin{enumerate}
  \def\labelenumi{(\alph{enumi})}
  \setcounter{enumi}{2}
  \tightlist
  \item
    How many total subsets of features are there (including empty set)?
    Hint: Each feature is either in or out = \(2^{10}\)
  \end{enumerate}
\end{itemize}

\textbf{43.} A neural network architect is designing a network with: -
Choice of 3 activation functions - Choice of 4 optimization algorithms -
Number of hidden layers: 1, 2, 3, or 4 - Learning rate: 5 options

How many different configurations are possible?

\begin{center}\rule{0.5\linewidth}{0.5pt}\end{center}

** Science Connection: Genetics**

\textbf{44.} DNA has 4 bases: A, T, G, C. A codon is a sequence of 3
bases.

\begin{itemize}
\item
  \begin{enumerate}
  \def\labelenumi{(\alph{enumi})}
  \tightlist
  \item
    How many different codons are possible?
  \end{enumerate}
\item
  \begin{enumerate}
  \def\labelenumi{(\alph{enumi})}
  \setcounter{enumi}{1}
  \tightlist
  \item
    There are 20 amino acids. Why does this work out nicely?
  \end{enumerate}
\end{itemize}

\textbf{45.} In a genetics experiment, a researcher needs to test all
pairwise combinations of 8 genes. - (a) How many pairs must be tested? -
(b) If each test takes 2 hours, how long will the full experiment take?

\begin{center}\rule{0.5\linewidth}{0.5pt}\end{center}

\hypertarget{pascals-triangle-connection}{%
\subsubsection{Pascal's Triangle
Connection}\label{pascals-triangle-connection}}

\textbf{46.} Row 5 of Pascal's Triangle is: 1, 5, 10, 10, 5, 1

These are
\(\binom{5}{0}, \binom{5}{1}, \binom{5}{2}, \binom{5}{3}, \binom{5}{4}, \binom{5}{5}\)

Verify: \(\binom{5}{2} = 10\) using the formula.

\textbf{47.} The sum of row \(n\) of Pascal's Triangle is \(2^n\).

\begin{itemize}
\item
  \begin{enumerate}
  \def\labelenumi{(\alph{enumi})}
  \tightlist
  \item
    Verify for row 5: \(1 + 5 + 10 + 10 + 5 + 1 = ?\)
  \end{enumerate}
\item
  \begin{enumerate}
  \def\labelenumi{(\alph{enumi})}
  \setcounter{enumi}{1}
  \tightlist
  \item
    This represents all possible subsets of a 5-element set. Why?
  \end{enumerate}
\end{itemize}

\begin{center}\rule{0.5\linewidth}{0.5pt}\end{center}

\hypertarget{part-e-putting-it-together-ml-applications-10-problems}{%
\subsection{Part E: Putting It Together --- ML Applications (10
problems)}\label{part-e-putting-it-together-ml-applications-10-problems}}

** Accuracy, Precision, Recall**

\textbf{48.} A model classifying tumors as ``malignant'' or ``benign''
produces:

\begin{longtable}[]{@{}lll@{}}
\toprule\noalign{}
& Predicted Malignant & Predicted Benign \\
\midrule\noalign{}
\endhead
\bottomrule\noalign{}
\endlastfoot
Actually Malignant & 85 (True Positive) & 15 (False Negative) \\
Actually Benign & 20 (False Positive) & 880 (True Negative) \\
\end{longtable}

Calculate: - (a) \textbf{Accuracy} = (TP + TN) / Total - (b)
\textbf{Precision} = TP / (TP + FP) --- ``Of predicted positives, how
many are correct?'' - (c) \textbf{Recall} = TP / (TP + FN) --- ``Of
actual positives, how many did we catch?'' - (d) In cancer detection, is
it worse to have low precision or low recall?

\textbf{49.} Another model has: - Accuracy: 95\% - Precision: 70\% -
Recall: 99\%

What does this tell you about the model's behavior?

\begin{center}\rule{0.5\linewidth}{0.5pt}\end{center}

** Expected Value**

\textbf{50.} Expected value: \(E[X] = \sum x_i \cdot P(x_i)\)

A game costs €5 to play. You roll a die: - Roll 6: win €20 - Roll 5: win
€10 - Roll anything else: win €0

\begin{itemize}
\item
  \begin{enumerate}
  \def\labelenumi{(\alph{enumi})}
  \tightlist
  \item
    Calculate the expected winnings.
  \end{enumerate}
\item
  \begin{enumerate}
  \def\labelenumi{(\alph{enumi})}
  \setcounter{enumi}{1}
  \tightlist
  \item
    What is the expected profit (winnings minus cost)?
  \end{enumerate}
\item
  \begin{enumerate}
  \def\labelenumi{(\alph{enumi})}
  \setcounter{enumi}{2}
  \tightlist
  \item
    Is this a fair game?
  \end{enumerate}
\end{itemize}

\textbf{51.} A trading algorithm has: - 60\% chance of €100 profit -
25\% chance of breaking even - 15\% chance of €200 loss

\begin{itemize}
\item
  \begin{enumerate}
  \def\labelenumi{(\alph{enumi})}
  \tightlist
  \item
    Calculate the expected value per trade.
  \end{enumerate}
\item
  \begin{enumerate}
  \def\labelenumi{(\alph{enumi})}
  \setcounter{enumi}{1}
  \tightlist
  \item
    Over 100 trades, what is the expected total profit?
  \end{enumerate}
\end{itemize}

\begin{center}\rule{0.5\linewidth}{0.5pt}\end{center}

** Bayes' Theorem Preview**

\textbf{52.} Bayes' Theorem: \[P(A|B) = \frac{P(B|A) \cdot P(A)}{P(B)}\]

A factory has two machines: - Machine A produces 60\% of items, with 2\%
defect rate - Machine B produces 40\% of items, with 5\% defect rate

An item is selected at random and found to be defective. What is the
probability it came from Machine A?

\begin{itemize}
\item
  \begin{enumerate}
  \def\labelenumi{(\alph{enumi})}
  \tightlist
  \item
    Find \(P(\text{defective})\) = P(def\textbar A)P(A) +
    P(def\textbar B)P(B)
  \end{enumerate}
\item
  \begin{enumerate}
  \def\labelenumi{(\alph{enumi})}
  \setcounter{enumi}{1}
  \tightlist
  \item
    Apply Bayes' Theorem
  \end{enumerate}
\end{itemize}

\begin{center}\rule{0.5\linewidth}{0.5pt}\end{center}

** Scientific Sampling**

\textbf{53.} A biologist wants to estimate the number of fish in a lake
using capture-recapture: - Day 1: Capture 100 fish, tag them, release -
Day 2: Capture 80 fish, find 10 are tagged

Using the proportion:
\(\frac{\text{tagged in sample}}{\text{sample size}} \approx \frac{\text{total tagged}}{\text{population}}\)

Estimate the total fish population.

\begin{center}\rule{0.5\linewidth}{0.5pt}\end{center}

** Cross-Validation**

\textbf{54.} In 5-fold cross-validation, a model's accuracy on each fold
is:

Fold 1: 88\%, Fold 2: 85\%, Fold 3: 90\%, Fold 4: 87\%, Fold 5: 85\%

\begin{itemize}
\item
  \begin{enumerate}
  \def\labelenumi{(\alph{enumi})}
  \tightlist
  \item
    What is the mean accuracy?
  \end{enumerate}
\item
  \begin{enumerate}
  \def\labelenumi{(\alph{enumi})}
  \setcounter{enumi}{1}
  \tightlist
  \item
    What is the standard deviation?
  \end{enumerate}
\item
  \begin{enumerate}
  \def\labelenumi{(\alph{enumi})}
  \setcounter{enumi}{2}
  \tightlist
  \item
    Report the result as ``mean ± SD''
  \end{enumerate}
\end{itemize}

\begin{center}\rule{0.5\linewidth}{0.5pt}\end{center}

** A/B Testing**

\textbf{55.} A website runs an A/B test on a new checkout button: -
Version A (old): 1000 visitors, 50 purchases (5\% conversion) - Version
B (new): 1000 visitors, 65 purchases (6.5\% conversion)

\begin{itemize}
\item
  \begin{enumerate}
  \def\labelenumi{(\alph{enumi})}
  \tightlist
  \item
    What is the relative improvement of B over A?
  \end{enumerate}
\item
  \begin{enumerate}
  \def\labelenumi{(\alph{enumi})}
  \setcounter{enumi}{1}
  \tightlist
  \item
    Is this difference ``real'' or could it be random chance? (This
    requires statistical tests we'll learn later, but intuitively, is 15
    extra purchases out of 1000 a big deal?)
  \end{enumerate}
\end{itemize}

\begin{center}\rule{0.5\linewidth}{0.5pt}\end{center}

\end{document}
